% mage-auto-heading: the-canonical-move-permitted-alternatives
% mage-auto-heading: compare-the-decision-with-the-facts
% mage-auto-heading: preserve-the-cases-a-decision-must-distinguish
% Authoritative LaTeX manuscript; Markdown is a generated compatibility projection.

% <!-- part-title: Modeling -->

% <!-- chapter-title: Decision Models: What Is Allowed? -->

\chapter{Decision Models: What Is Allowed?}

The preceding models already expressed obligations: a user interface must reach its database through the service layer, and a job must pass through \texttt{RUNNING} before \texttt{DONE}. A decision model focuses directly on the alternatives an engineer permits, requires, or forbids. It can answer which service may call another, or whether a user is entitled to upload a document.

\section{The Canonical Move: Permitted Alternatives}\label{heading:the-canonical-move-permitted-alternatives}

% <!-- point: permission-has-many-canonical-forms-unified-by-mood-not-shape | Decision tables, rule systems, access-control matrices, policy graphs, feature models, and constraint systems can express which alternatives are permitted; their shape does not determine that purpose. -->

Engineering represents permitted alternatives in several familiar forms: decision tables that map conditions to allowed actions, rule systems, access-control matrices that record which subject may act on which object,\citep{lampson1974} policy graphs, feature models that constrain which combinations of options are valid, and constraint systems that declare which assignments are admissible. These representations differ in shape, but each can make the rules for a decision explicit.

% <!-- point: same-arrow-different-mood-is-the-sharpest-case-of-topology-not-fixing-meaning | An observed-call edge and a permitted-call edge can look identical while answering different questions; in the illustrated models, a missing edge means no recorded call or no permitted call respectively. -->

The meaning of an edge therefore depends on the model\textquotesingle s purpose. In an observed call graph, \texttt{A~→~B} records that A calls B; in a policy graph, the same arrow can mean that A \emph{may} call B. These graphs let engineers compare what happens with what is permitted. For the models illustrated in \autoref{canonical-decision}, a missing edge means that no such call is recorded in the observed graph, or that no such call is permitted by the policy. This reading depends on the model\textquotesingle s declared scope and rules; the arrow\textquotesingle s shape alone does not supply it.

\begin{figure}[htbp]
\centering
\includegraphics[width=\linewidth]{book/latex/assets/canonical-decision-47ef23e866.pdf}
\caption{\emph{The same arrow can record observation or permission.} In the illustrated scope, a missing edge means no recorded call in the observed graph or no permitted call in the policy graph. The model's declared meaning determines which interpretation applies.}
\label{canonical-decision}
\end{figure}

% <!-- point: a-declared-policy-turns-is-this-allowed-into-an-evaluation | Once permission is declared, the claims a policy designer wants become expressible — permitted/forbidden, mutual exclusion, prerequisites, consistency, completeness, least privilege — and "is this allowed?" becomes an evaluation rather than an argument. -->

Once permission is declared, the claims a policy designer wants become expressible: which alternatives are permitted or forbidden, which are mutually exclusive, which imply prerequisites, whether the policy is internally consistent, whether it is complete over the cases that can arise, and whether it achieves least privilege. The analyses follow --- evaluating a rule against a case, checking or solving constraints, detecting conflicts between rules, testing completeness and consistency. A declared policy turns "is this allowed?" into an evaluation rather than an argument.

\section{Compare the Decision with the Facts}\label{heading:compare-the-decision-with-the-facts}

Suppose a report service may call a renderer, while the public web interface may call only the report service. A policy graph records those permitted pairs. An observed call from the web interface directly to the renderer would then be a candidate violation: the observation supplies the fact, and the policy supplies the rule. The comparison also depends on the observation\textquotesingle s coverage. A missing observed call does not prove that the call can never occur. \autoref{observed-vs-declared} separates the observed call from the policy used to judge it.

\begin{figure}[htbp]
\centering
\includegraphics[width=\linewidth]{book/latex/assets/observed-vs-declared-4186370076.pdf}
\caption{\emph{Observation and intent answer different questions.} The observed system can reveal that A calls C; the decision model can state whether that relationship is permitted. Agreement establishes correspondence, not correctness.}
\label{observed-vs-declared}
\end{figure}

A policy can agree with every observed call and still permit more than the design should allow. Correspondence asks whether the observations satisfy the declared policy. Choosing the policy asks whether those permissions meet the engineering requirements. Both questions matter, and neither representation answers the other automatically.

\section{Preserve the Cases a Decision Must Distinguish}\label{heading:preserve-the-cases-a-decision-must-distinguish}

The input representation can make a decision wrong before any rule is evaluated. Return to the prepaid and unlimited subscriptions introduced earlier. For a simple policy with one applicable entitlement, \autoref{simple-admission-policy} states when to admit a request and whether to charge finite credit.

{\footnotesize
\begin{longtable}[]{@{}
  >{\raggedright\arraybackslash}p{(\linewidth - 4\tabcolsep) * \real{0.3333}}
  >{\raggedright\arraybackslash}p{(\linewidth - 4\tabcolsep) * \real{0.3333}}
  >{\raggedright\arraybackslash}p{(\linewidth - 4\tabcolsep) * \real{0.3333}}@{}}
\caption[Admission under one entitlement]{\emph{Admission under one entitlement.} Unlimited entitlement admits covered work without drawing down finite credit; a finite balance must cover the requested work.}\label{simple-admission-policy}\tabularnewline
\toprule\noalign{}
\begin{minipage}[b]{\linewidth}\raggedright
Entitlement
\end{minipage} & \begin{minipage}[b]{\linewidth}\raggedright
Requested work
\end{minipage} & \begin{minipage}[b]{\linewidth}\raggedright
Decision
\end{minipage} \\
\midrule\noalign{}
\endhead
\bottomrule\noalign{}
\endlastfoot
Unlimited & Any amount covered by the subscription & Admit without drawing down finite credit \\
Finite balance & Enough credit for the request & Admit and charge the required credit \\
Finite balance & Insufficient credit for the request & Reject for insufficient credit \\
\end{longtable}
}

The table needs the entitlement kind as well as the balance. If a previous step passes only a number, zero cannot distinguish an exhausted finite allowance from unlimited use. The decision code cannot recover a distinction that its input has discarded. DocAble supplies two worked applications: a declared service-call policy and an upload-admission failure that required a richer decision model.

% <!-- box-family: inset -->

% <!-- inset-domain: WORKED EXAMPLE -->

% <!-- box-weight: deep-dive -->

% mage-box-prefix:

\begin{bookbox}{inset}{Worked example --- Permission at two system boundaries}
\label{inset-docable-decisions}
DocAble accepts documents for accessibility remediation, sends work among internal services, and returns corrected artifacts. Two decisions arise at different boundaries: which service may call another, and whether a user\textquotesingle s entitlement permits a requested upload. Both can be modeled as permission, but they need different representations.

\textbf{Which services may communicate?}

DocAble declares which of its services may communicate with which. The model is a graph whose edges carry the normative reading: an edge means the call is part of the sanctioned policy, and an absent edge means the policy does not permit it.

\textbf{Service-flow model · Decision}

\begin{itemize}
\tightlist
\item
  \textbf{Engineering question} --- Which service may call which, and what may be reached from outside?
\end{itemize}

\begin{itemize}
\tightlist
\item
  \textbf{Model} --- a declared graph of services and the permitted edges between them; the deployment is wired from the graph rather than the graph inferred from the deployment.
\end{itemize}

\begin{itemize}
\tightlist
\item
  \textbf{Property} --- internal communication is permitted only along a declared edge; one declared node is the public entry.
\end{itemize}

\begin{itemize}
\tightlist
\item
  \textbf{Quality attribute} --- least privilege; architectural integrity; security.
\end{itemize}

% <!-- point: the-service-flow-model-is-a-declared-source-of-truth-not-a-firewall | The declared graph is the single source of truth for permitted inter-service edges, with one node declared publicly reachable and every other edge authenticated; a completeness check requires each gated edge to be fully specified, and the deployment config is generated from the graph. | terms: soft-vs-hard-governance -->

\autoref{service-flow-model} draws the permitted calls. One node is declared publicly reachable; every other edge is an authenticated internal call. The model is the single source of truth for those edges, and the deployment is generated from it rather than inferred from it --- so a completeness check can require every gated edge to be fully specified before the deployment is wired. That check is real and enforced, but it governs the \emph{declaration\textquotesingle s} internal consistency, not the running traffic: the model declares the topology; it is not a runtime firewall.

\begin{figure}[htbp]
\centering
\includegraphics[width=\linewidth]{book/latex/assets/service-flow-model-7634fa9071.pdf}
\caption{\emph{Service-flow model.} A declared graph of which service may reach which. The public web tier is the only node facing the outside; an absent edge means the relationship is not permitted by the model, and enforcement is separate.}
\label{service-flow-model}
\end{figure}

% <!-- point: observed-versus-declared-is-correspondence-not-correctness | The observed call graph reveals that A calls C; the decision model states whether A may call C; agreement is correspondence, not correctness — a faithfully-obeyed rule can still be the wrong rule. -->

The declared graph supplies the policy against which an observed call graph could be checked.

% <!-- point: docable-does-not-reconcile-observed-traffic-so-correspondence-is-the-open-question | DocAble generates config from the declared graph but does not reconcile observed traffic against it at runtime; whether the deployed services stay on the permitted edges, and what a departure should cost, is Alignment's question. -->

Agreement between observation and policy would establish \emph{correspondence} --- the implementation matching the declared intent --- not \emph{correctness}: a faithfully-obeyed rule can still be the wrong rule. DocAble generates its deployment from the declared graph, but nothing at runtime reconciles observed traffic back against the declaration. Whether the deployed services stay on the declared edges, and what a departure should cost, is Alignment\textquotesingle s question in \magevalue{chapter:alignment}.

\textbf{Which entitlement authorizes an upload?}

% <!-- point: the-entitlement-model-was-born-from-a-production-failure | The service-flow policy was declared in advance; DocAble's upload-admission decision acquired its model in the opposite order, after a production failure exposed a representation that could no longer express the decision. -->

The service-flow policy states an architectural decision. A different decision, whether a user may upload a document, exposed the cost of losing a domain distinction as the product changed.

% <!-- point: upload-entitlement-flows-through-groups | A user receives upload entitlement through groups — finite-credit or unlimited — so the admission decision depends on two facts: how many credits remain and whether any entitlement is unlimited. -->

DocAble admits an upload only when the user is entitled to it. A user receives that entitlement through groups. A group may carry a finite credit balance, drawn down as pages are processed, or it may provide unlimited use. The admission decision therefore depends on two different facts: how many finite credits are available, and whether any applicable entitlement is unlimited (\autoref{entitlement-domain}).

\begin{figure}[htbp]
\centering
\includegraphics[width=\linewidth]{book/latex/assets/entitlement-domain-14ff87d0d6.pdf}
\caption{\emph{The upload-entitlement domain.} A user belongs to one or more groups; a group carries either a finite credit balance or unlimited use. The admission decision needs both facts, so a representation that can express only the finite balance cannot answer the question the decision asks.}
\label{entitlement-domain}
\end{figure}

% <!-- point: the-scalar-was-adequate-until-the-domain-changed | DocAble originally had only finite-credit groups, so a scalar sum was an adequate admission representation; when unlimited groups arrived, no explicit model of the admission decision existed to carry the change. -->

DocAble originally represented upload entitlement in finite credits alone. In that system, reducing a user\textquotesingle s groups to a single number --- the sum of their credit balances --- was adequate for deciding whether a document could be processed. The product later gained unlimited groups, and the domain change altered the admission question: entitlement was no longer a single finite quantity. No explicit model of the upload-admission decision existed at the time. The decision lived in several implementation sites --- a server-side reduction, a client-side preflight, the server\textquotesingle s admission gate --- each with its own hand-rolled representation of the user\textquotesingle s groups.

% <!-- point: the-old-reduction-observed-the-distinction-and-discarded-it | The legacy reduction read each group's unlimited flag and used it to exclude unlimited groups from the credit sum; its scalar output had no place for "unlimited," so an unlimited user became zero credits — a representation failure, not missing information. -->

One of those representations did not evolve. In August 2026, a user whose one group carried unlimited entitlement was blocked from uploading a thirty-eight-page document by a paywall: \emph{the upload needs 38 credits, and you have 0}. The root cause was a legacy server function that reduced a user\textquotesingle s groups to one scalar, the sum of their finite credit balances. The function was not ignorant of unlimited entitlement. It read each group\textquotesingle s unlimited flag twice, using it once to \emph{exclude} unlimited groups from the sum and once to compute a display label. Its output type simply had no place to put the fact, so the distinction died at the return statement. A purely-unlimited user became \texttt{0} credits; the client compared 0 against 38 and rejected the upload; the server\textquotesingle s authoritative admission logic, which handled unlimited correctly, was never reached. This was a representation failure, not an information-availability failure. The implementation held the fact and threw it away.

% <!-- point: the-production-failure-produced-the-model | The root-cause analysis made the admission question explicit and produced an explicit entitlement decision model — with property tests, a cross-runtime parity corpus, and a shape check — in roughly five hours from RCA to deletion of the lossy reduction. -->

The root-cause analysis made the admission question explicit: given the user\textquotesingle s group entitlements and the requested work, should this upload be admitted, and under which entitlement? We introduced an entitlement decision model to answer it. Property tests checked the model, a cross-runtime parity corpus supplied shared test cases to check that server and client followed the same decision table, and a further check validated the model\textquotesingle s shape. We deleted the legacy scalar reduction without a compatibility shim. The display code then derived its values from the model, and the admission gates used the model\textquotesingle s decision. From root-cause analysis to deletion of the lossy representation took roughly five hours.

\textbf{Entitlement decision model · Decision}

\begin{itemize}
\tightlist
\item
  \textbf{Engineering question} --- May this user upload this document, and under which entitlement?
\end{itemize}

\begin{itemize}
\tightlist
\item
  \textbf{Model} --- per-group entitlements that keep finite and unlimited as distinct cases; the output is a structured admission decision: admitted or not, the group that authorizes it, and the ordered plan by which the verdict was reached.
\end{itemize}

\begin{itemize}
\tightlist
\item
  \textbf{Property} --- an unlimited entitlement can never be reduced to zero credits; the displayed balance is a projection of the model, not an authority for admission.
\end{itemize}

\begin{itemize}
\tightlist
\item
  \textbf{Quality attribute} --- decision correctness; fidelity between the domain and its representations.
\end{itemize}

% <!-- point: the-repair-was-model-to-model | The repair changed the representation through which the decision was made, not merely the failing conditional: the revised model preserves the finite/unlimited distinction and states the admission decision explicitly, and the implementation was then changed to consume it. -->

The repair did more than patch the conditional that failed. It changed the representation through which the decision was made. Upload entitlement could no longer be represented as a scalar credit balance: the revised model preserves finite and unlimited entitlement as distinct cases and represents the admission decision explicitly --- whether the upload is admitted, which group supplies the entitlement, and the ordered plan used to reach the verdict. The implementation was then changed to consume that model. The same triggering state that once read as \texttt{0} credits now reads as \emph{admit by unlimited entitlement}. The finite-credit scalar still exists where that quantity is the actual question, displaying remaining credits, but it is no longer used as a model of upload permission (\autoref{entitlement-model-evolution}).

\begin{figure}[htbp]
\centering
\includegraphics[width=\linewidth]{book/latex/assets/entitlement-model-evolution-31ef2274cf.pdf}
\caption{\emph{Upload admission must distinguish finite and unlimited entitlement.} The old reduction passed only a finite balance to the client, so an unlimited user was rejected. The revised model records the entitlement kind and returns an admission decision.}
\label{entitlement-model-evolution}
\end{figure}

% <!-- point: the-model-was-part-of-the-repair | The failure revealed a decision made through a representation that could no longer express the system's domain; making the decision explicit forced the representation to preserve the distinction the system had acquired. -->

The revised decision model was part of the repair, rather than merely documentation of it. The failure revealed that an engineering decision was being made through a representation that could no longer express the system\textquotesingle s domain, and making the decision explicit forced the representation to preserve the distinction the system had acquired. The opening section defined a model as a purposeful reduction that preserves the distinctions its question requires. The old scalar was a reasonable model of \emph{how many finite credits remain}. Once unlimited entitlement was introduced, that scalar was no longer a sufficient model of \emph{may this user upload}. The gap between those two questions then became a production failure.

% <!-- point: an-explicit-model-would-have-made-the-obligation-visible | The author changed the domain in one place and did not carry the change into every representation of the same decision; when the domain contains a consequential distinction, every representation used to reason about that decision must be capable of expressing it. -->

A note in the first person: I changed the domain in one place --- groups could now be unlimited --- and did not carry the change into every representation of the same decision. An explicit model would not have made the change automatically correct; it would have made the inconsistency visible as an engineering obligation. When the domain contains a consequential distinction, every representation used to reason about that decision must be capable of expressing it. \magevalue{chapter:alignment} gives the recurring pattern here --- a failure converted into explicit structure that later work inherits --- a name.

The two models give a decision an explicit domain: permitted service pairs in the first case, distinct entitlement alternatives in the second. That makes the decision inspectable and checkable. It also exposes what must change when the intended policy or the product\textquotesingle s domain changes.
\end{bookbox}

Decision models delimit the permitted alternatives while making the facts required for the decision explicit. The final section steps back from one representation to ask how a family of models can answer questions together.
